\section{Outcome (e): Addition and multiplication rules}

\subsection{What the outcome asks}

The exam tests two rules and the habit of choosing the right one. The addition rule
(inclusion-exclusion) gives the probability that at least one of several events occurs; the exam
states it for two or three events, usually with one quantity missing, and asks for that quantity or
for a derived one such as ``neither'' or ``exactly one''. The multiplication rule (chain rule)
gives the probability that several events occur in sequence, each conditioned on the ones before it;
the exam uses it for draws without replacement and for claim histories over successive years.
Many questions combine the two: split the outcome into disjoint cases with the addition rule, then
compute each case with the multiplication rule. A tree diagram is the bookkeeping device for the
multiplication rule, and its branches are the cases that the addition rule adds up.

\subsection{Addition rule}

\begin{keyfact}[General addition rule, two events]
For any events $A$ and $B$,
\[
\Prob(A \cup B) = \Prob(A) + \Prob(B) - \Prob(A \cap B).
\]
The subtraction removes the double count of outcomes that lie in both events. The rule holds
whether or not $A$ and $B$ are independent or disjoint.
\end{keyfact}

\begin{keyfact}[General addition rule, three events]
For any events $A$, $B$, $C$,
\[
\Prob(A \cup B \cup C) = \Prob(A) + \Prob(B) + \Prob(C)
 - \Prob(A \cap B) - \Prob(A \cap C) - \Prob(B \cap C)
 + \Prob(A \cap B \cap C).
\]
Sign pattern: singles added, pairs subtracted, triple added. Seven terms on the right; the exam
typically gives six of them (or the union) and asks for the seventh.
\end{keyfact}

Special cases reduce the number of terms.

\begin{itemize}[leftmargin=2em]
  \item \textbf{Disjoint (mutually exclusive).} If $A \cap B = \emptyset$ then
  $\Prob(A \cup B) = \Prob(A) + \Prob(B)$. For three pairwise disjoint events the union is the sum
  of the three probabilities.
  \item \textbf{Independent.} If $A$ and $B$ are independent then $\Prob(A \cap B) = \Prob(A)\Prob(B)$, so
  $\Prob(A \cup B) = \Prob(A) + \Prob(B) - \Prob(A)\Prob(B)$.
  \item \textbf{Complement shortcut.} By De Morgan, $\cc{(A \cup B)} = \cc{A} \cap \cc{B}$, so
  $\Prob(A \cup B) = 1 - \Prob(\cc{A} \cap \cc{B})$ and $\Prob(\text{neither}) = 1 - \Prob(A \cup B)$.
  For $n$ events, $\Prob(\text{at least one}) = 1 - \Prob(\text{none})$; when the events are
  independent this is $1 - \prod_{i=1}^{n}(1 - \Prob(A_i))$.
  \item \textbf{Bound.} Without any information about intersections,
  $\Prob(A_1 \cup \cdots \cup A_n) \le \sum_{i=1}^{n} \Prob(A_i)$ (Boole's inequality), with
  equality only when the events are pairwise disjoint.
\end{itemize}

\begin{keyfact}[Exactly one of two events]
\[
\Prob(\text{exactly one of } A, B) = \Prob(A) + \Prob(B) - 2\,\Prob(A \cap B)
 = \Prob(A \cup B) - \Prob(A \cap B).
\]
Given any two of $\Prob(A)+\Prob(B)$, $\Prob(A\cap B)$, $\Prob(A\cup B)$,
$\Prob(\text{exactly one})$, the remaining ones follow.
\end{keyfact}

\begin{keyfact}[Counting identity for three events]
Write $e_1$, $e_2$, $e_3$ for the probabilities that exactly one, exactly two, exactly three of
$A, B, C$ occur. Then
\[
\Prob(A \cup B \cup C) = e_1 + e_2 + e_3, \qquad
\Prob(A) + \Prob(B) + \Prob(C) = e_1 + 2e_2 + 3e_3.
\]
The second identity holds because an outcome in exactly $k$ of the events is counted $k$ times in
the sum of the marginals. The exam uses these two linear equations to recover $e_2$ or $e_3$ from
``at least one'' and ``exactly one'' data.
\end{keyfact}

\paragraph{Solving-for-unknown patterns.} The exam rarely asks for the union directly. The
standard forms are:

\begin{enumerate}[leftmargin=2em]
  \item Given $\Prob(A \cup B)$, $\Prob(A)$, $\Prob(B)$: the intersection is
  $\Prob(A \cap B) = \Prob(A) + \Prob(B) - \Prob(A \cup B)$.
  \item Given $\Prob(\cc{A} \cap \cc{B})$ (``neither''): first $\Prob(A \cup B) = 1 - \Prob(\cc{A} \cap \cc{B})$,
  then proceed as in pattern 1.
  \item Given $\Prob(\text{exactly one})$ and the two marginals: $\Prob(A \cap B) =
  \tfrac12\bigl[\Prob(A) + \Prob(B) - \Prob(\text{exactly one})\bigr]$.
  \item Three events with six of the seven inclusion-exclusion terms and either the union or
  ``none'' given: solve the seven-term equation for the missing term.
  \item Three events with ``at least one'' and ``exactly one'' given together with the sum of the
  marginals: solve the two linear equations in $e_2$ and $e_3$.
\end{enumerate}

\subsection{Multiplication rule}

\begin{keyfact}[General multiplication rule, two events]
For any events $A$ and $B$ with positive probability,
\[
\Prob(A \cap B) = \Prob(A)\,\Prob(B \mid A) = \Prob(B)\,\Prob(A \mid B).
\]
This is the definition of conditional probability rearranged. Use whichever ordering matches
the information given.
\end{keyfact}

\begin{keyfact}[Chain rule, three or more events]
\[
\Prob(A_1 \cap A_2 \cap A_3) = \Prob(A_1)\,\Prob(A_2 \mid A_1)\,\Prob(A_3 \mid A_1 \cap A_2),
\]
and in general
$\Prob(A_1 \cap \cdots \cap A_n) = \Prob(A_1) \prod_{k=2}^{n} \Prob(A_k \mid A_1 \cap \cdots \cap A_{k-1})$.
Each factor is conditioned on everything that happened before it.
\end{keyfact}

\begin{keyfact}[Independence collapses the chain]
If $A_1, \dots, A_n$ are mutually independent, every conditional factor equals the marginal and
$\Prob(A_1 \cap \cdots \cap A_n) = \Prob(A_1)\Prob(A_2)\cdots\Prob(A_n)$.
The product of marginals is a special case of the chain rule, not a separate rule. Multiply
unconditional probabilities only when the problem states independence or the setup implies it
(sampling with replacement, separate policyholders described as independent).
\end{keyfact}

\paragraph{Sequential sampling without replacement.} When items are drawn one at a time without
replacement, the composition of the pool changes after each draw, so the draws are dependent. The
chain rule handles this directly: the $k$th factor is the number of favourable items remaining
divided by the number of items remaining. From a pool of $N$ items of which $m$ are of type $X$,
\[
\Prob(\text{first two both type } X) = \frac{m}{N}\cdot\frac{m-1}{N-1}, \qquad
\Prob(\text{first is } X,\ \text{second is not}) = \frac{m}{N}\cdot\frac{N-m}{N-1}.
\]

\begin{keyfact}[Symmetry of draws without replacement]
The unconditional probability that the $k$th draw is of type $X$ equals the probability that the
first draw is of type $X$, namely $m/N$, for every $k \le N$. Every ordering of the $N$ items is
equally likely, so the item in position $k$ is a uniformly random item from the pool. Once the
results of earlier draws are known, the conditional probability changes.
\end{keyfact}

\paragraph{Tree diagrams.} A tree diagram lists the first event's outcomes as branches, then each
second event's outcomes as sub-branches, and so on. Each branch is labelled with a conditional
probability given the path that leads to it; the probability of a complete path is the product of
its branch labels, which is exactly the chain rule. Because the paths are disjoint, the probability
of any event is the sum of the probabilities of the paths in it, which is the addition rule for
disjoint cases. As a table, the tree for a two-year claim history with
$\Prob(C_1) = 0.20$, $\Prob(C_2 \mid C_1) = 0.35$, $\Prob(C_2 \mid \cc{C_1}) = 0.15$ has one row per path:

\begin{center}
\begin{tabular}{@{}llrr@{}}
\toprule
Year 1 (branch) & Year 2 (branch) & Product & Path probability \\
\midrule
Claim ($0.20$)    & Claim ($0.35$)    & $0.20 \times 0.35$ & $0.070$ \\
Claim ($0.20$)    & No claim ($0.65$) & $0.20 \times 0.65$ & $0.130$ \\
No claim ($0.80$) & Claim ($0.15$)    & $0.80 \times 0.15$ & $0.120$ \\
No claim ($0.80$) & No claim ($0.85$) & $0.80 \times 0.85$ & $0.680$ \\
\midrule
 & & Total & $1.000$ \\
\bottomrule
\end{tabular}
\end{center}

The path probabilities sum to $1.000$, which is the standard check. Any event is a set of rows and
its probability is the sum of those rows.

\subsection{Combining the rules}

The most common exam structure is a two-stage experiment: a classification or a first draw, then a
second outcome whose probability depends on the first. The method is fixed: list the disjoint cases
for the first stage (the branches of the tree), which must cover every possibility without
overlapping; for each case compute $\Prob(\text{case}) \times \Prob(\text{target} \mid \text{case})$
with the multiplication rule; add the results with the addition rule for disjoint events.
Written out for cases $B_1, \dots, B_n$ that partition the sample space and a target event $A$,
\[
\Prob(A) = \sum_{i=1}^{n} \Prob(B_i)\,\Prob(A \mid B_i).
\]
This is the law of total probability. Outcome (g) develops it fully and adds Bayes' theorem for
reversing the conditioning; for this outcome it is enough to recognise that ``split into cases,
multiply along each path, add the paths'' is the same computation. Example: 70\% of policyholders
are low risk with claim probability $0.10$ and 30\% are high risk with claim probability $0.30$;
the probability that a randomly selected policyholder has a claim is
$0.70 \times 0.10 + 0.30 \times 0.30 = 0.07 + 0.09 = 0.16$.

A question that gives conditional probabilities along a chain and asks for ``at least one'' over
the chain is usually fastest through the ``none'' path: multiply along it and subtract from 1.

\subsection{Worked examples}

\begin{example}[Three coverages, solve for the triple intersection]
An insurer's policyholders may hold auto, homeowners, and life coverage. Among them, 60\% hold
auto, 45\% hold homeowners, 30\% hold life, 25\% hold both auto and homeowners, 15\% hold both auto
and life, 10\% hold both homeowners and life, and 5\% hold none of the three.
Calculate the percentage of policyholders who hold all three coverages.
\par\smallskip\noindent (A)~10\% \quad (B)~5\% \quad (C)~8\% \quad (D)~12\% \quad (E)~15\%
\end{example}
\begin{solution}
Let $A$, $H$, $L$ be the events that a policyholder holds auto, homeowners, life. The
complement of the union is ``none'', so $\Prob(A \cup H \cup L) = 1 - 0.05 = 0.95$. Three-event
inclusion-exclusion with $x = \Prob(A \cap H \cap L)$:
\[
0.95 = 0.60 + 0.45 + 0.30 - 0.25 - 0.15 - 0.10 + x = 0.85 + x,
\]
so $x = 0.10$. Answer \textbf{(A)}.
\end{solution}

\begin{example}[Claims in two successive years]
For a policyholder, the probability of at least one claim in year 1 is $0.20$. If the
policyholder has a claim in year 1, the probability of a claim in year 2 is $0.35$; if there is
no claim in year 1, the probability of a claim in year 2 is $0.15$.
Calculate the probability that the policyholder has claims in both years or in neither year.
\par\smallskip\noindent (A)~0.68 \quad (B)~0.70 \quad (C)~0.80 \quad (D)~0.75 \quad (E)~0.93
\end{example}
\begin{solution}
Let $C_1$, $C_2$ be the claim events. By the multiplication rule,
\begin{align*}
\Prob(C_1 \cap C_2) &= \Prob(C_1)\Prob(C_2 \mid C_1) = 0.20 \times 0.35 = 0.07,\\
\Prob(\cc{C_1} \cap \cc{C_2}) &= \Prob(\cc{C_1})\Prob(\cc{C_2} \mid \cc{C_1}) = 0.80 \times 0.85 = 0.68.
\end{align*}
``Both'' and ``neither'' are disjoint, so the addition rule gives $0.07 + 0.68 = 0.75$.
Answer \textbf{(D)}. As a check, ``exactly one year'' is $0.20 \times 0.65 + 0.80 \times 0.15 = 0.25$
and $0.75 + 0.25 = 1$.
\end{solution}

\begin{example}[Two policies of the same type]
A portfolio contains 20 policies: 10 auto, 6 homeowners, and 4 life. An auditor selects two
policies at random without replacement.
Calculate the probability that the two selected policies are of the same type.
\par\smallskip\noindent (A)~0.30 \quad (B)~0.35 \quad (C)~0.33 \quad (D)~0.38 \quad (E)~0.40
\end{example}
\begin{solution}
``Same type'' splits into three disjoint cases: both auto, both homeowners, both life. Each case
is a chain-rule product, and the cases are added:
\[
\frac{10}{20}\cdot\frac{9}{19} + \frac{6}{20}\cdot\frac{5}{19} + \frac{4}{20}\cdot\frac{3}{19}
 = \frac{90 + 30 + 12}{380} = \frac{132}{380} = \frac{33}{95} \approx 0.3474.
\]
Answer \textbf{(B)}. Counting unordered pairs gives the same value, $(45 + 15 + 6)/190$.
\end{solution}

\begin{example}[Type of the second policy drawn]
From the portfolio of the previous example (10 auto, 6 homeowners, 4 life), the auditor selects two
policies at random without replacement.
Calculate the probability that the second policy selected is a life policy.
\par\smallskip\noindent (A)~0.15 \quad (B)~0.16 \quad (C)~0.19 \quad (D)~0.21 \quad (E)~0.20
\end{example}
\begin{solution}
Condition on the first draw. Either the first policy is life (probability $4/20$, leaving 3 life
among 19) or it is not (probability $16/20$, leaving 4 life among 19):
\[
\Prob(\text{second is life}) = \frac{4}{20}\cdot\frac{3}{19} + \frac{16}{20}\cdot\frac{4}{19}
 = \frac{12 + 64}{380} = \frac{76}{380} = \frac{4}{20} = 0.20.
\]
Answer \textbf{(E)}. The result equals the first-draw probability $4/20$. This is not a coincidence:
if nothing is known about the first draw, the second policy is a uniformly random policy from the
20, so its type distribution is the same as the pool's. Algebraically,
$\frac{m}{N}\cdot\frac{m-1}{N-1} + \frac{N-m}{N}\cdot\frac{m}{N-1} = \frac{m(N-1)}{N(N-1)} = \frac{m}{N}$.
Recognising the symmetry saves the computation; the computation serves as a check.
\end{solution}

\begin{example}[Three coverages from ``at least one'' and ``exactly one'']
A survey of an insurer's policyholders finds that 60\% hold auto coverage, 35\% hold homeowners
coverage, and 20\% hold life coverage. Also, 80\% hold at least one of the three coverages and
50\% hold exactly one.
Calculate the percentage of policyholders who hold all three coverages.
\par\smallskip\noindent (A)~0\% \quad (B)~10\% \quad (C)~5\% \quad (D)~15\% \quad (E)~25\%
\end{example}
\begin{solution}
Let $e_1, e_2, e_3$ be the proportions holding exactly one, two, three coverages. ``At least one''
gives $e_1 + e_2 + e_3 = 0.80$ and $e_1 = 0.50$, so $e_2 + e_3 = 0.30$. The sum of the marginals
counts each policyholder once per coverage held:
\[
e_1 + 2e_2 + 3e_3 = 0.60 + 0.35 + 0.20 = 1.15
 \quad\Longrightarrow\quad 2e_2 + 3e_3 = 0.65.
\]
Subtract twice the first equation from the second: $e_3 = 0.65 - 2(0.30) = 0.05$, and then
$e_2 = 0.25$. Answer \textbf{(C)}. The individual pairwise intersections are neither needed nor
determined by the data; only their sum $e_2 + 3e_3 = 0.40$ is.
\end{solution}

\begin{example}[Complement of a union from count data]
Of 200 liability claims closed last year, 90 involved property damage, 60 involved bodily injury,
and 120 involved property damage or bodily injury or both. A claim is selected at random.
Calculate the probability that the selected claim involved neither property damage nor bodily
injury.
\par\smallskip\noindent (A)~0.40 \quad (B)~0.15 \quad (C)~0.25 \quad (D)~0.45 \quad (E)~0.60
\end{example}
\begin{solution}
Work in counts, then convert. The union has 120 claims, so its complement has $200 - 120 = 80$
claims and the probability is $80/200 = 0.40$. Answer \textbf{(A)}. The intersection is
$90 + 60 - 120 = 30$ claims, and ``exactly one'' is $120 - 30 = 90$ claims. The 90 and 60 are
counts, not percentages; treating them as proportions produces a nonsense answer.
\end{solution}

\begin{example}[Chain rule over three years]
A policyholder's claim history is modelled year by year. The probability of no claim in year 1 is
$0.80$. Given no claim in year 1, the probability of no claim in year 2 is $0.85$. Given no claim
in years 1 and 2, the probability of no claim in year 3 is $0.90$.
Calculate the probability that the policyholder has at least one claim during the three years.
\par\smallskip\noindent (A)~0.29 \quad (B)~0.35 \quad (C)~0.45 \quad (D)~0.39 \quad (E)~0.55
\end{example}
\begin{solution}
Let $N_k$ be the event of no claim in year $k$. By the chain rule,
\[
\Prob(N_1 \cap N_2 \cap N_3) = \Prob(N_1)\Prob(N_2 \mid N_1)\Prob(N_3 \mid N_1 \cap N_2)
 = 0.80 \times 0.85 \times 0.90 = 0.612.
\]
At least one claim is the complement: $1 - 0.612 = 0.388$. Answer \textbf{(D)}. The problem gives no
conditional probabilities after a claim, so the paths containing a claim cannot be computed and
the complement route is the only one available.
\end{solution}

\subsection{Traps}

\begin{trap}[Adding without subtracting the intersection]
$\Prob(A \cup B) = \Prob(A) + \Prob(B)$ only when $A$ and $B$ are disjoint. Policyholders can hold
two coverages and a claim can involve two causes. Unless the
question says ``mutually exclusive'' or the events cannot occur together by construction, subtract
$\Prob(A \cap B)$. A quick check: if $\Prob(A) + \Prob(B) > 1$, the events cannot be disjoint.
\end{trap}

\begin{trap}[Multiplying marginals without independence]
$\Prob(A \cap B) = \Prob(A)\Prob(B)$ requires independence. Draws without replacement, successive
years of the same policyholder's claims, and risk class followed by claim status are all dependent.
The correct factor for the second event is the conditional probability given the first. If the
question supplies $\Prob(B \mid A)$, it is telling you the events are not independent.
\end{trap}

\begin{trap}[Confusing the joint with the conditional]
$\Prob(A \cap B)$ is the probability that both happen; $\Prob(B \mid A)$ is the probability of $B$
in the world where $A$ has already happened. ``The probability that a policyholder with a claim
in year 1 has a claim in year 2'' is $\Prob(C_2 \mid C_1)$. ``The probability that a policyholder
has claims in both years'' is $\Prob(C_1 \cap C_2) = \Prob(C_1)\Prob(C_2 \mid C_1)$. The
conditional is never smaller than the joint.
\end{trap}

\begin{trap}[Sign pattern in three-event inclusion-exclusion]
Singles are added, pairwise intersections are subtracted, and the triple intersection is added
back. The common errors are subtracting the triple intersection or omitting it. Write the full
seven-term equation before substituting. Check the result: each pairwise intersection must be at
least the triple intersection, and each marginal at least any pairwise intersection containing it.
\end{trap}

\begin{trap}[Counts, percentages, and the base]
Some questions give raw counts (90 of 200 claims), some give percentages of the whole group, and
some give percentages of a subgroup (``of those with auto coverage, 40\% also hold homeowners'').
The third kind is a conditional probability, not a joint probability, and must be multiplied by
the subgroup's proportion before it enters the addition rule. With counts, work in counts and
divide by the total once at the end.
\end{trap}

\subsection{Practice problems}

\begin{practice}
An insurer's records show that 55\% of policyholders are male, 40\% are over age 50, and 20\% are
males over age 50. A policyholder is selected at random.
Calculate the probability that the selected policyholder is a female who is age 50 or younger.
\par\smallskip\noindent (A)~0.05 \quad (B)~0.25 \quad (C)~0.20 \quad (D)~0.30 \quad (E)~0.45
\end{practice}

\begin{practice}
Among an insurer's policyholders, 50\% hold auto coverage, 40\% hold homeowners coverage, 30\%
hold life coverage, 20\% hold auto and homeowners, 15\% hold auto and life, 10\% hold homeowners
and life, and 5\% hold all three.
Calculate the percentage of policyholders who hold none of the three coverages.
\par\smallskip\noindent (A)~10\% \quad (B)~15\% \quad (C)~25\% \quad (D)~30\% \quad (E)~20\%
\end{practice}

\begin{practice}
A claims adjuster has 12 open files: 5 auto, 4 homeowners, and 3 life. The adjuster selects three
files at random without replacement to review first.
Calculate the probability that none of the three selected files is a life file.
\par\smallskip\noindent (A)~0.30 \quad (B)~0.35 \quad (C)~0.38 \quad (D)~0.42 \quad (E)~0.45
\end{practice}

\begin{practice}
A batch of 25 newly issued policies contains 7 with a data-entry error. A quality reviewer examines
the policies one at a time in random order without replacement.
Calculate the probability that the third policy examined contains an error.
\par\smallskip\noindent (A)~0.28 \quad (B)~0.24 \quad (C)~0.26 \quad (D)~0.30 \quad (E)~0.32
\end{practice}

\begin{practice}
For a policyholder, the probability of a claim in year 1 is $0.15$. In any year, if the
policyholder had a claim in the preceding year, the probability of a claim is $0.40$; if the
policyholder had no claim in the preceding year, the probability of a claim is $0.10$.
Calculate the probability that the policyholder has a claim in exactly one of years 1 and 2.
\par\smallskip\noindent (A)~0.135 \quad (B)~0.150 \quad (C)~0.200 \quad (D)~0.175 \quad (E)~0.235
\end{practice}

\begin{practice}
Among an auto insurer's policyholders, 45\% carry collision coverage, 35\% carry comprehensive
coverage, and 50\% carry exactly one of the two.
Calculate the percentage of policyholders who carry both collision and comprehensive coverage.
\par\smallskip\noindent (A)~5\% \quad (B)~15\% \quad (C)~10\% \quad (D)~20\% \quad (E)~30\%
\end{practice}

\begin{practice}
A life insurer classifies applicants as standard, preferred, or substandard; 60\% of applicants are
standard, 30\% are preferred, and 10\% are substandard. The probability of a claim within the
first policy year is $0.08$ for a standard applicant, $0.03$ for a preferred applicant, and $0.20$
for a substandard applicant.
Calculate the probability that a randomly selected applicant has a claim within the first policy
year.
\par\smallskip\noindent (A)~0.058 \quad (B)~0.065 \quad (C)~0.085 \quad (D)~0.103 \quad (E)~0.077
\end{practice}

\subsection{Answers to practice problems}

\begin{enumerate}[leftmargin=2em]
  \item \textbf{(B)} $0.25$. $\Prob(M \cup O) = 0.55 + 0.40 - 0.20 = 0.75$; female and 50 or younger
  is the complement of the union, $1 - 0.75 = 0.25$.
  \item \textbf{(E)} $20\%$. Union $= 0.50 + 0.40 + 0.30 - 0.20 - 0.15 - 0.10 + 0.05 = 0.80$; none
  $= 1 - 0.80 = 0.20$.
  \item \textbf{(C)} $0.38$. Nine non-life files among 12: $\frac{9}{12}\cdot\frac{8}{11}\cdot\frac{7}{10}
  = \frac{504}{1320} = \frac{21}{55} \approx 0.3818$.
  \item \textbf{(A)} $0.28$. By symmetry of draws without replacement, the third policy is a uniformly
  random policy from the batch, so the probability is $7/25 = 0.28$.
  \item \textbf{(D)} $0.175$. Claim then no claim: $0.15 \times 0.60 = 0.090$; no claim then claim:
  $0.85 \times 0.10 = 0.085$; disjoint, so add: $0.175$.
  \item \textbf{(B)} $15\%$. Exactly one $= \Prob(A) + \Prob(B) - 2\Prob(A \cap B)$, so
  $0.50 = 0.80 - 2x$ and $x = 0.15$.
  \item \textbf{(E)} $0.077$. Split by class and multiply along each path:
  $0.60(0.08) + 0.30(0.03) + 0.10(0.20) = 0.048 + 0.009 + 0.020 = 0.077$.
\end{enumerate}

\subsection{One-page summary}

\begin{itemize}[leftmargin=2em]
  \item Two events: $\Prob(A \cup B) = \Prob(A) + \Prob(B) - \Prob(A \cap B)$. Always subtract
  the intersection unless the events are disjoint.
  \item Three events: add the three singles, subtract the three pairs, add the triple. Seven terms;
  the exam hides one.
  \item Complement of the union is ``neither'' or ``none'':
  $\Prob(\cc{A} \cap \cc{B}) = 1 - \Prob(A \cup B)$. ``At least one'' $= 1 - \Prob(\text{none})$.
  \item Exactly one of two events: $\Prob(A) + \Prob(B) - 2\Prob(A \cap B) = \Prob(A \cup B) - \Prob(A \cap B)$.
  \item Three events with $e_1, e_2, e_3$ (exactly one, two, three): union $= e_1 + e_2 + e_3$ and
  sum of marginals $= e_1 + 2e_2 + 3e_3$.
  \item Multiplication rule: $\Prob(A \cap B) = \Prob(A)\Prob(B \mid A) = \Prob(B)\Prob(A \mid B)$.
  Chain rule: $\Prob(A_1 \cap A_2 \cap A_3) = \Prob(A_1)\Prob(A_2 \mid A_1)\Prob(A_3 \mid A_1 \cap A_2)$;
  each factor conditions on all earlier events.
  \item Under independence every conditional factor becomes the marginal and the chain is a plain
  product of marginals. Do not use the plain product otherwise.
  \item Without replacement: $k$th factor is (favourable items remaining)/(items remaining).
  Unconditionally, the $k$th draw has the same type distribution as the first: $m/N$.
  \item Tree diagram: branch labels are conditional probabilities; a path's probability is the
  product of its labels; an event's probability is the sum of its paths. Paths sum to 1.
  \item Two-stage problems: split into disjoint cases, multiply along each case, add. This is the
  law of total probability, $\Prob(A) = \sum_i \Prob(B_i)\Prob(A \mid B_i)$ (outcome (g)).
  \item Joint versus conditional: $\Prob(A \cap B) = \Prob(A)\Prob(B \mid A) \le \Prob(B \mid A)$.
  A percentage ``of those with $A$'' is a conditional probability. With counts, divide by the
  total once at the end.
\end{itemize}
